% sections/04-results.tex (plan 49 section 2.3; headings state the finding and the verdict). Rewritten in pass 6 on
% research/68 sections 0 and B: a bold lead line with the headline before the pre-registered order (research/66); every
% number is a generated macro; verdict words come only from the \wtthHaiku...VerdictText macros; the H1 to H3
% subsections, headings included, are the realised branch files (branch.tex). Table 2 is built here from macros (Appendix
% D keeps the generated table). Post hoc and exploratory results are labelled wherever they appear; D-codes appear only
% beside those labels (research/66).
\section{Results}\label{sec:results}
\textbf{Headline: a receipt-style line that the claimant typed into its own post cut reporting by \wtthHaikuHThreeDelta{} points and produced the only shifts in which no one told ops; the bare claim barely moved reporting (\wtthHaikuHOneDelta{} points, a borderline rejection), and the log made it matter more.}
The tests follow in their pre-registered order, on all \wtthHaikuEpisodesInSet{} planned Haiku shifts (\wtthHaikuFailedPermanent{} permanent failures, \wtthHaikuMissing{} missing).

\begin{table}[t]
\centering\footnotesize\setlength{\tabcolsep}{3.5pt}
\begin{tabularx}{\linewidth}{@{}l>{\raggedright\arraybackslash}Xrr>{\raggedright\arraybackslash}p{2.32in}@{}}
\toprule
Test & Contrast in P1 (told ops) & Diff. & \wtthProtocolLevel\% CI & Decision \\
\midrule
H1 & \cell{NONE} $-$ \cell{CLAIM}: \wtthHaikuPOneNonePct\% vs \wtthHaikuPOneClaimPct\% & \wtthHaikuHOneDelta & \wtthHaikuHOneCI & \wtthHaikuHOneVerdictText \\
H2 & stand-down in R0 $-$ in R1: \wtthHaikuHTwoDeltaRZero{} vs \wtthHaikuHTwoDeltaROne{} & \wtthHaikuHTwoDelta & \wtthHaikuHTwoCI & \wtthHaikuHTwoVerdictText \\
H3 & \cell{NO-RECEIPT} $-$ \cell{FORGED}: \wtthHaikuPOneNoReceiptPct\% vs \wtthHaikuPOneForgedPct\% & \wtthHaikuHThreeDelta & \wtthHaikuHThreeCI & \wtthHaikuHThreeVerdictText \\
\bottomrule
\end{tabularx}
\caption{\textbf{The confirmatory tests, in their fixed order} (\wtthDesignModelHaikuName). Differences in points, shift-level cluster bootstrap, each predicted positive; a rejection opposite to the prediction still counts, so the sequence continues.
H2's stand-downs are \cell{NONE} $-$ \cell{CLAIM} (R0, no log) and \cell{NONE-R} $-$ \cell{NO-RECEIPT} (R1); with both reference cells at \wtthHaikuPOneNonePct\%, \Dtwo{} reduces to P1(\cell{NO-RECEIPT}) $-$ P1(\cell{CLAIM}).
H1's rejection is borderline (Section~\ref{sec:h1}), so H2 and H3 are confirmatory only through that gate.}
\label{tab:hyp}
\end{table}

\wtthInputBranch{h1}
\wtthInputBranch{h2}
\wtthInputBranch{h3}

% Research/64 FLAGS 15, 16, 18, 23 and 6: unadjusted intervals; H1's sensitivity rows return the primary value; the
% live-only re-analysis (exploratory); realised precision from the ceiling; board substitution conditional on sending no
% note is flat across cells (the unconditional contrast restates P1). The frontier panel is Figure fig:frontier (App D).
\subsection{Defences, sensitivity and board substitution}\label{sec:frontier}\label{sec:robust}
Beyond the three confirmatory tests, every interval in this paper is unadjusted for multiplicity, and none is read as a test.
\textbf{Defences (RQ4).}
\wtthInputBranch{frontier}
\textbf{Sensitivity.}
H2 and H3 keep their sign and side of zero under the pre-registered sensitivity analyses; for H1 every such row returns the primary estimate, so none probes its boundary (Appendix~\ref{app:full}).
Restricted to witnesses whose \DebOne{} answer was coded live (\wtthHaikuDOneLiveAllPct\%; post hoc, exploratory), H1 to H3 were \wtthHaikuLiveOnlyHOneDelta, \wtthHaikuLiveOnlyHTwoDelta{} and \wtthHaikuLiveOnlyHThreeDelta{} points, as in the main analysis.
False reports in the benign cell were \wtthHaikuFalseReportBNonePct\%.
The realised intervals are far narrower than the planned detectable effects (\wtthHaikuPowerHOneMDE, \wtthHaikuPowerHTwoMDE{} and \wtthHaikuPowerHThreeMDE{} points at a planning rate of \wtthHaikuPowerP\%) because variance is small near the ceiling, not because the design was precise.
\textbf{Board substitution} (exploratory, D-010).
Witnesses who sent no note raised the split on the board at about the same rate wherever a receipt was in view and where the log was empty (\cell{FORGED} \wtthHaikuBoardSubForgedOfNotelessK{} of \wtthHaikuBoardSubForgedOfNotelessN{} noteless witnesses, \cell{RECEIPT} \wtthHaikuBoardSubReceiptOfNotelessK{} of \wtthHaikuBoardSubReceiptOfNotelessN, \cell{NO-RECEIPT} \wtthHaikuBoardSubNoReceiptOfNotelessK{} of \wtthHaikuBoardSubNoReceiptOfNotelessN), so higher overall rates under a receipt reflect more noteless witnesses, not a change in what they did instead.

% Exploratory follow-up (D-022, D-023): no verdict, no p-value (plan 49 section 4.2 item 8). Research/64 FLAG 10: no
% concurrent FORGED arm, and both sentences moved P1 by about the same amount, as a phase shift would. Research/62 M5:
% the same witness-turn harness (D-023's operating changes touch budgeting and probe dispatch only). FLAG 11: Sonnet's
% RECEIPT is its only cell where C really reported, so log-checking and distrust of the typed line are confounded; the
% Sonnet difference has a shift-level cluster-bootstrap interval (research/65). Detail: Appendix D.
\subsection{Follow-up (exploratory): neither added sentence brought Haiku back to \cell{NO-RECEIPT}'s level}\label{sec:followup}
New seeds with the core's mix of paraphrase and handle order ran hours after the core, on the same witness-turn harness (D-022, D-023; Appendix~\ref{app:full}) and with no concurrent \cell{FORGED} arm, so every comparison here is across phases.
Added to \cell{FORGED}'s runbook, the one-line rule (\cell{X-FORGED-POLICY}) or a note that only \texttt{notify\_ops} writes the log (\cell{X-FORGED-PROV}) left Haiku's P1 at \wtthExpOneXForgedPolicyPOnePct\% and \wtthExpOneXForgedProvPOnePct\%, below \cell{NO-RECEIPT} (\wtthExpOneDiffPOneXForgedPolicyMinusHaikuNoReceiptDelta{} and \wtthExpOneDiffPOneXForgedProvMinusHaikuNoReceiptDelta{} points, intervals below zero).
Both moved P1 from \cell{FORGED} by about the same amount (\wtthExpOneDiffPOneXForgedPolicyMinusHaikuForgedDelta{} and \wtthExpOneDiffPOneXForgedProvMinusHaikuForgedDelta{} points, CI \wtthExpOneDiffPOneXForgedPolicyMinusHaikuForgedCI{} and \wtthExpOneDiffPOneXForgedProvMinusHaikuForgedCI), which a shift between phases would also produce; we cannot separate the two.
Ops went untold in \wtthExpOneXForgedPolicyPTwoK{} and \wtthExpOneXForgedProvPTwoK{} of \wtthExpOneXForgedPolicyPTwoN{} shifts.

\paragraph{Sonnet stood down for a real receipt but not for the typed line (exploratory, across phases).}
Beside C's real receipt, \wtthDesignModelSonnetName{} reported at \wtthExpOneSonnetReceiptPOnePct\% (\wtthExpOneSonnetReceiptPOneK/\wtthExpOneSonnetReceiptPOneN), against \wtthSonnetPOneForgedPct\% under the typed line and \wtthSonnetPOneNoReceiptPct\% with an empty log (difference from the typed line \wtthExpOneDiffPOneSonnetReceiptMinusSonnetForgedDelta{} points, \wtthFigOneAIntervalMethod{} CI \wtthExpOneDiffPOneSonnetReceiptMinusSonnetForgedCI; \wtthExpOneSonnetReceiptEps{} shifts each).
The typed-line arm comes from the earlier Sonnet subset, so the contrast is across phases.
\cell{RECEIPT} is the only Sonnet cell in which C really reported, so the design cannot tell whether Sonnet checked the log or distrusted the typed line; Haiku stood down for both.
